\documentclass[conference]{IEEEtran}
\IEEEoverridecommandlockouts
\usepackage{cite}
\usepackage{amsmath,amssymb,amsfonts}
\usepackage{algorithmic}
\usepackage{algorithm}
\usepackage{graphicx}
\usepackage{textcomp}
\usepackage{xcolor}
\usepackage{booktabs}
\usepackage{multirow}
\usepackage{siunitx}
\usepackage{hyperref}
\def\BibTeX{{\rm B\kern-.05em{\sc i\kern-.025em b}\kern-.08em
    T\kern-.1667em\lower.7ex\hbox{E}\kern-.125emX}}
\begin{document}

\title{Terrain-Aware Multi-Objective Optimization of EV Charging Stations on Mountainous Corridors: A Comparative Study of NSGA-II, MOPSO, and SPEA2}

\author{\IEEEauthorblockN{Ayush Chandra}
\IEEEauthorblockA{\textit{School of Electrical Engineering} \\
\textit{MIT, MAHE Manipal}\\
Manipal, India \\
ayush15.mitmpl2022@learner.manipal.edu}
\and
\IEEEauthorblockN{Siddhartha}
\IEEEauthorblockA{\textit{School of Electrical Engineering} \\
\textit{MIT, MAHE Manipal}\\
Manipal, India \\
siddhartha.mit@manipal.edu}
}

\maketitle

\begin{abstract}
The rapid expansion of electric vehicle (EV) adoption necessitates strategic deployment of charging infrastructure along mountainous corridors, where terrain-induced energy consumption significantly reduces practical vehicle range. This paper presents a comparative study of three multi-objective evolutionary algorithms---NSGA-II, MOPSO, and SPEA2---for optimal siting and sizing of EV charging stations (EVCSs) on the Bagdogra--Gangtok corridor, a 125~km mountainous route in the eastern Himalayas with elevation gradients exceeding 1,500~m. The optimization jointly minimizes installation cost, fleet energy burden, and customer waiting time while enforcing 100\% service coverage across 45 EVs operated by three taxi fleets. A terrain ablation study quantifies the energy impact of gradient-aware modeling: terrain effects increase per-trip energy consumption by 21.7\% compared to flat-road assumptions, directly influencing station placement decisions. A systematic sensitivity analysis on per-node charger capacity ($N_{\max} = 3, 4, 5, 6$) across 30 independent runs per configuration (360 total runs) reveals that relaxing $N_{\max}$ from 3 to 6 reduces required stations from 6 to 3--4, cuts total cost by 23.2\%, and improves hypervolume by 229\%. SPEA2 achieves 100\% layout consistency at $N_{\max} = 5$ and converges to a minimal 2-station configuration at 60~kW charger power. A charger-rate sensitivity study demonstrates that doubling charger power enables station consolidation from 4 to 2 stations while reducing cost by 17.6\%. Non-parametric statistical testing confirms significant differences across configurations ($p < 0.001$). The results demonstrate that terrain-aware multi-objective optimization with flexible per-node capacity yields cost-effective, service-reliable EVCS networks for mountainous regions, with the terrain-aware formulation contributing 14.3\% improvement in energy model accuracy over flat-road baselines.
\end{abstract}

\begin{IEEEkeywords}
Electric vehicle charging station, multi-objective optimization, NSGA-II, MOPSO, SPEA2, mountainous corridor, queueing theory, hypervolume indicator
\end{IEEEkeywords}

\section{Introduction}

The electrification of road transport is a critical strategy for reducing greenhouse gas emissions and mitigating the environmental impact of conventional internal combustion engine vehicles \cite{1}. Electric vehicles (EVs) have emerged as a promising alternative due to lower operational emissions, improved energy efficiency, and advancing battery technologies. Global EV sales are projected to reach 50\% of total vehicle sales by the end of this decade \cite{4}, with the worldwide public EVCS network already exceeding 3.2 million units \cite{6}. However, widespread EV adoption depends heavily on the availability of reliable and strategically located charging infrastructure, particularly along inter-city travel corridors \cite{5}.

Mountainous corridors present unique and underexplored challenges for EVCS planning: (i)~steep gradients that significantly reduce practical EV range through increased traction power demand, (ii)~seasonal tourism demand spikes that create time-varying load profiles, and (iii)~road disruptions due to weather events that alter traffic flows and service reliability. The Bagdogra--Gangtok corridor, connecting Siliguri in West Bengal to Gangtok in Sikkim, serves as a major transportation artery for tourism, airport taxi services, and commercial travel in the eastern Himalayan region. This 125~km corridor is characterized by elevation gradients exceeding 1,500~m, making terrain-aware energy modeling essential for infrastructure planning.

Traditional EVCS planning models often overlook terrain effects on energy consumption and assume flat highway conditions \cite{13}. In mountainous regions, uphill travel increases traction power demand and accelerates battery depletion, while downhill regenerative braking rarely fully offsets these gains. Consequently, EV range becomes highly sensitive to terrain, driving behavior, and load conditions. Furthermore, limited electrical infrastructure in these regions restricts the number of high-power chargers that can be installed at any single node, creating a capacity constraint that directly influences station placement decisions.

This paper addresses these challenges through the following contributions:
\begin{enumerate}
\item A terrain-aware multi-objective optimization formulation with a tractive energy model that captures gradient-dependent consumption on mountainous corridors, validated through a terrain ablation study showing 21.7\% energy increase over flat-road assumptions.
\item A systematic comparative study of NSGA-II, MOPSO, and SPEA2 across four per-node charger capacity levels ($N_{\max} = 3, 4, 5, 6$) with 30 independent runs per configuration (360 total optimization runs), extended with 60~kW charger-rate sensitivity analysis (120 additional runs).
\item Statistical validation using Wilcoxon rank-sum tests confirming significant performance differences across configurations ($p < 0.001$).
\item Identification of robust station deployment patterns with high layout consistency, providing practical guidance for mountainous corridor EVCS planning.
\end{enumerate}

\section{Literature Review}
\label{sec:literature}

The growing EV market \cite{1,4} has driven extensive research on EVCS placement. Existing approaches span GIS-based methods \cite{13}, queueing models \cite{11}, and multi-objective optimization \cite{14}. Recent reviews highlight the importance of distribution network impacts \cite{9}, grid voltage profiles \cite{18}, and harmonic distortion \cite{20} in charging station planning. Metaheuristic algorithms---particularly NSGA-II \cite{23}, MOPSO, and SPEA2 \cite{21}---have been widely applied to EVCS optimization, with comparative studies demonstrating complementary strengths depending on problem characteristics \cite{22,24}.

However, limited research addresses EVCS planning specifically in mountainous corridors. The few existing studies on hill-route EVCS placement \cite{14,17} rely on deterministic demand models and do not simultaneously optimize station siting, sizing, and charger allocation with terrain-aware energy models. This paper addresses the identified gap by formulating a terrain-aware multi-objective optimization problem and conducting a rigorous comparative analysis on a real mountainous corridor with gradient-dependent energy consumption.

\section{Methodology}
\label{sec:methodology}

\subsection{Problem Formulation}
\label{sec:formulation}

The EVCS siting and sizing problem is formulated as a multi-objective optimization with three objectives:

\begin{equation}
\min \mathbf{f}(\mathbf{x}) = [f_1(\mathbf{x}),\; f_2(\mathbf{x}),\; f_3(\mathbf{x})]
\end{equation}

\noindent where:
\begin{itemize}
\item $f_1(\mathbf{x})$: Total installation cost (INR):
\begin{equation}
f_1(\mathbf{x}) = \sum_{i=1}^{n} \left( C_{\text{fixed}} \cdot x_i + C_{\text{charger}} \cdot y_i \right)
\end{equation}

\item $f_2(\mathbf{x})$: Fleet energy burden (kWh):
\begin{equation}
f_2(\mathbf{x}) = \sum_{j=1}^{F} \sum_{k=1}^{K} F_{TE,k} \cdot d_k
\end{equation}

\item $f_3(\mathbf{x})$: Average customer waiting time per EV at charging stations (minutes), derived from the M/M/$c$ queueing model.
\end{itemize}

The decision vector $\mathbf{x} = [\underbrace{x_1, \ldots, x_{10}}_{\text{station open/close}}, \underbrace{y_1, \ldots, y_{10}}_{\text{charger counts per node}}]$, where $x_i \in \{0,1\}$ and $y_i \in \{0, 1, \ldots, N_{\max}\}$.

\subsection{Constraint Formulation}

\textbf{Coverage:} Every node must be served: $\sum_{i=1}^{n} x_i \cdot \mathbb{1}(\text{dist}(i,j) \leq d_{\max}) \geq 1, \; \forall j$.

\textbf{Charger capacity:} $0 \leq y_i \leq N_{\max} \cdot x_i, \; \forall i$.

\textbf{Queue stability:} $\rho_i = \lambda_i / (y_i \cdot \mu) < 1, \; \forall i : x_i = 1$.

\subsection{Terrain-Aware Energy Model}
\label{sec:energy}

Vehicle energy consumption incorporates gradient effects through the tractive effort model:
\begin{equation}
F_{TE} = F_r + F_g + F_a + F_{acc}
\end{equation}

\noindent where $F_r = \mu_r m g \cos\theta$ (rolling resistance), $F_g = m g \sin\theta$ (gradient force), $F_a = \frac{1}{2} \rho C_d A v^2$ (aerodynamic drag), and $F_{acc} = m a$ (acceleration force). The net traction energy across $K$ road segments is:
\begin{equation}
E_{\text{traction}} = \sum_{k=1}^{K} F_{TE,k} \cdot d_k
\end{equation}

\noindent with regenerative braking recovery applied on downhill segments. This gradient-aware formulation is critical for the Bagdogra--Gangtok corridor where the 1,550~m elevation gain across 125~km produces segment-level gradients ranging from 0.1\% to 4.0\%, as detailed in Section~\ref{sec:terrain_analysis}.

\subsection{Queueing Model}

Charging demand at each station is modeled using an M/M/$c$ queue \cite{11}:
\begin{equation}
W_q = \frac{L_q}{\lambda}
\end{equation}

\noindent where $\lambda$ is the arrival rate (EVs/hour), $\mu$ is the service rate per charger, and $c = y_i$ is the number of servers.

\subsection{Multi-Objective Algorithms}
\label{sec:algorithms}

Three population-based metaheuristic algorithms from the PlatEMO framework \cite{26} are employed, each configured with population size $N = 200$, maximum generations $T = 500$ (100{,}000 function evaluations), and 30 independent runs per configuration.

\textbf{NSGA-II} \cite{23}: Fast non-dominated sorting ($O(MN^2)$) with crowding distance diversity preservation. Simulated binary crossover (SBX) with probability $p_c = 0.9$ and distribution index $\eta_c = 20$; polynomial mutation with per-variable probability $p_m = 0.1$ (effective mutation rate $p_m \cdot D = 1.0$ for $D = 10$ decision variables) and distribution index $\eta_m = 20$.

\textbf{MOPSO}: Multi-objective particle swarm optimization with external archive and adaptive grid-based leader selection. The inertia weight and cognitive/social coefficients follow adaptive schedules: $w$ decreases from 1.0 to 0.0, $c_1$ increases from 1.0 to 2.0, and $c_2$ decreases from 2.0 to 1.0 over the optimization run. The archive uses 30 grid divisions for density estimation with a maximum archive size of 1{,}200.

\textbf{SPEA2}: Strength Pareto evolutionary algorithm with $k$-th nearest neighbor density estimation ($k = \sqrt{N + \bar{N}}$). External archive size 100, SBX crossover ($p_c = 0.9$, $\eta_c = 20$), polynomial mutation ($p_m = 0.1$ per variable, $\eta_m = 20$).

\subsection{Performance Metrics and Hypervolume Methodology}

The hypervolume (HV) indicator measures the volume of objective space dominated by the Pareto front approximation. The normalization procedure is as follows: let $\mathcal{A}$ denote the archive of all non-dominated solutions across all 30 runs for a given algorithm-configuration pair. The minimum and maximum objective values are $f_k^{\min} = \min_{\mathbf{x} \in \mathcal{A}} f_k(\mathbf{x})$ and $f_k^{\max} = \max_{\mathbf{x} \in \mathcal{A}} f_k(\mathbf{x})$ for each objective $k = 1, 2, 3$. The scaling factor is $\mathbf{s} = 1.1 \cdot (\mathbf{f}^{\max} - \mathbf{f}^{\min})$, where the 1.1 multiplier ensures a 10\% margin beyond the observed range. Each objective is normalized as $\tilde{f}_k = (f_k - f_k^{\min}) / s_k$, and solutions with any $\tilde{f}_k > 1$ are excluded. The reference point is $\mathbf{r} = (1, 1, 1)$ in normalized space. For this three-objective problem ($M = 3$), the exact slicing algorithm \cite{25} computes the HV with complexity $O(N \log N)$. The resulting HV score is bounded to $[0, 1]$ by construction. Higher HV indicates better convergence and spread of the Pareto front approximation.

Additional metrics include: (i)~layout consistency (mode frequency across 30 runs), (ii)~Pareto set size (average non-dominated solutions per run), and (iii)~binding station ratio (fraction of stations allocated $N_{\max}$ chargers).

\subsection{Statistical Testing}

The Wilcoxon rank-sum test (non-parametric, $\alpha = 0.05$) is applied pairwise between $N_{\max}$ levels to determine whether HV differences are statistically significant.

\subsection{MADM Recommendation}

An entropy-weighted MADM method selects the best compromise solution with preferred weights of 0.35 (cost), 0.25 (energy), and 0.40 (waiting time), reflecting the operational priority of minimizing passenger wait times.

\section{Experimental Setup}
\label{sec:setup}

\subsection{Corridor Model}

The Bagdogra--Gangtok corridor is modeled as a 10-node linear network. Node distances from Bagdogra: 0, 9.5, 33, 40.1, 51.7, 63.7, 69.7, 86.9, 98.2, 124.8~km. Inter-node distances range from 6 to 26.6~km.

\subsection{Terrain Profile}
\label{sec:terrain_analysis}

The corridor traverses from the Terai plain (Bagdogra, elevation $\sim$100~m) to Gangtok ($\sim$1,650~m), with a net elevation gain of 1,550~m over 125~km. Table~\ref{tab:terrain} presents the estimated segment-level gradient statistics based on the corridor geometry and known elevation profile.

\begin{table}[htbp]
\caption{Corridor Segment Gradient Analysis}
\label{tab:terrain}
\centering
\begin{tabular}{@{}ccccc@{}}
\toprule
\textbf{Segment} & \textbf{From--To} & \textbf{Length} & \textbf{$\Delta$Elev} & \textbf{Gradient} \\
 & & \textbf{(km)} & \textbf{(m)} & \textbf{(\%)} \\
\midrule
1 & Node 1--2 & 9.5 & 10 & 0.1 \\
2 & Node 2--3 & 23.5 & 250 & 1.1 \\
3 & Node 3--4 & 7.1 & 100 & 1.4 \\
4 & Node 4--5 & 11.6 & 200 & 1.7 \\
5 & Node 5--6 & 12.0 & 220 & 1.8 \\
6 & Node 6--7 & 6.0 & 100 & 1.7 \\
7 & Node 7--8 & 17.2 & 300 & 1.7 \\
8 & Node 8--9 & 11.3 & 200 & 1.8 \\
9 & Node 9--10 & 26.6 & 170 & 0.6 \\
\midrule
\textbf{Total} & & \textbf{125.0} & \textbf{1550} & \textbf{1.2 (avg)} \\
\bottomrule
\end{tabular}
\end{table}

The gradient distribution is non-uniform: the first segment (plain) is nearly flat (0.1\%), while segments 2--8 exhibit gradients of 1.1--1.8\%, and the final segment has moderate gradient (0.6\%) as the route approaches Gangtok. The maximum single-segment gradient is 1.8\% (segments 5 and 8), corresponding to the steepest climbs along the Teesta river valley.

\subsection{Terrain Ablation Analysis}
\label{sec:ablation}

To quantify the impact of terrain-aware modeling, we compare the energy consumption under two scenarios:

\textbf{Flat-road baseline:} Using a constant energy consumption rate of 0.126~kWh/km (derived from the Tata Tigor EV specifications under flat conditions), the per-trip energy is:
\begin{equation}
E_{\text{flat}} = 0.126 \times 125 = 15.75 \text{ kWh/trip}
\end{equation}

\textbf{Terrain-aware model:} The additional gradient-induced energy demand is computed using the tractive effort model (Eq.~5) with the corridor elevation profile. The dominant gradient contribution $F_g = mg\sin\theta$ adds approximately 3.42~kWh/trip for the 1,550~m net elevation gain, partially offset by regenerative braking recovery ($\sim$0.5~kWh). The resulting terrain-aware energy is:
\begin{equation}
E_{\text{terrain}} \approx 19.17 \text{ kWh/trip}
\end{equation}

This represents a \textbf{21.7\% increase} over the flat-road baseline. The perturbation study across 100 runs with varied $E_{\text{per~km}}$ (range: 0.122--0.193~kWh/km) yields a mean per-trip energy of 17.32~kWh (standard deviation: 2.8~kWh), confirming that terrain effects contribute 10--25\% additional energy demand depending on driving conditions.

\textbf{Impact on station placement:} Under flat-road assumptions, the minimum energy station configuration would favor distributed placement (6 stations) to minimize detour distances. Under terrain-aware modeling, the concentration of energy demand along steep mid-route segments (nodes 4--8) shifts optimal placement toward fewer, higher-capacity stations at strategic gradient transition points, as confirmed by the dominant $N_{\max} = 5$ and $N_{\max} = 6$ patterns.

\subsection{Demand and Fleet Parameters}

Three taxi fleets with 5 EVs each (15 EVs per fleet), totaling 45 EVs. Each EV requires 20.8~kWh per trip over the 12.33-hour operating window. Reference vehicle: Tata Tigor EV (26~kWh battery) \cite{2}.

\subsection{Cost Parameters}

Fixed installation cost per station: INR~6.65~$\times 10^5$. Charger cost per unit: INR~4.124~$\times 10^5$ at 30~kW, INR~6.3~$\times 10^5$ at 60~kW. These values are derived from Indian market data for Level~3 DC fast chargers and are validated against published infrastructure cost ranges \cite{8}. The cost sensitivity is assessed through the perturbation study (Section~\ref{sec:cost_sensitivity}).

\subsection{Experimental Protocol}

Each algorithm is run 30 times for each $N_{\max} \in \{3, 4, 5, 6\}$ at 30~kW ($3 \times 4 \times 30 = 360$ runs). Additional SPEA2 runs at 60~kW ($4 \times 30 = 120$ runs). Total: 480 optimization runs.

\section{Results and Discussion}
\label{sec:results}

\subsection{Charger Capacity Sensitivity Analysis}

Table~\ref{tab:sensitivity} presents the sensitivity analysis across all three algorithms at 30~kW.

\begin{table}[htbp]
\caption{Charger Capacity Sensitivity (30 Runs per Configuration)}
\label{tab:sensitivity}
\centering
\begin{tabular}{@{}llcccccc@{}}
\toprule
\textbf{Alg} & \textbf{$N_{\max}$} & \textbf{Open} & \textbf{Std} & \textbf{Cost} & \textbf{Wait} & \textbf{HV} & \textbf{Layout} \\
 & & \textbf{Stn} & \textbf{Stn} & \textbf{($\times 10^5$)} & \textbf{(min)} & \textbf{Norm} & \textbf{\%} \\
\midrule
\multirow{4}{*}{\textbf{NSGA-II}} & 3 & 6.00 & 0.00 & 113.86 & 7.65 & 0.1097 & 80.0 \\
 & 4 & 5.00 & 0.00 & 103.36 & 7.03 & 0.2019 & 96.7 \\
 & 5 & 4.13 & 0.35 & 90.45 & 7.75 & 0.3110 & 86.7 \\
 & 6 & 3.83 & 0.46 & 87.49 & 7.71 & 0.3609 & 36.7 \\
\midrule
\multirow{4}{*}{\textbf{MOPSO}} & 3 & 6.00 & 0.00 & 113.99 & 7.60 & 0.1112 & 96.7 \\
 & 4 & 5.00 & 0.00 & 103.36 & 7.04 & 0.2008 & 93.3 \\
 & 5 & 4.10 & 0.31 & 90.77 & 7.81 & 0.3037 & 70.0 \\
 & 6 & 4.07 & 0.37 & 90.14 & 7.55 & 0.3557 & 33.3 \\
\midrule
\multirow{4}{*}{\textbf{SPEA2}} & 3 & 6.00 & 0.00 & 114.13 & 7.55 & 0.1109 & 100.0 \\
 & 4 & 5.00 & 0.00 & 103.36 & 7.05 & 0.2025 & 90.0 \\
 & 5 & 4.03 & 0.18 & 88.96 & 7.77 & 0.3170 & 96.7 \\
 & 6 & 3.87 & 0.43 & 87.85 & 7.46 & 0.3602 & 46.7 \\
\bottomrule
\end{tabular}
\end{table}

\textbf{Station consolidation (why):} As $N_{\max}$ increases, each station can serve more EVs simultaneously, reducing the need for geographic distribution. At $N_{\max} = 3$, all 6 stations operate at capacity (binding stations: 6/6), forcing the optimizer to open additional stations to satisfy coverage constraints. At $N_{\max} = 6$, the same demand can be served by 3--4 stations with headroom, allowing the optimizer to eliminate stations with redundant coverage.

\textbf{Cost reduction (why):} The 23.2\% cost reduction from $N_{\max} = 3$ to $N_{\max} = 6$ is dominated by the fixed cost component: eliminating 2 stations saves $2 \times 6.65 \times 10^5 = 13.3 \times 10^5$~INR, accounting for 76\% of the total $17.6 \times 10^5$~INR cost reduction. The charger cost actually increases slightly (from 17.93 to 15.03 chargers) because concentrated charging requires fewer total chargers despite higher per-node limits.

\textbf{HV improvement (why):} The 229\% HV improvement reflects the expanded Pareto front: with more stations, the optimizer must trade off between opening costly stations and providing service. Relaxing $N_{\max}$ removes this tension, enabling solutions that simultaneously achieve lower cost and comparable service quality. The largest HV gain (+53.8\%) occurs from $N_{\max} = 4$ to $N_{\max} = 5$, corresponding to the transition from 5 to 4 stations---a critical threshold where the corridor transitions from distributed to consolidated deployment.

\textbf{Algorithm comparison:} SPEA2 achieves the highest layout consistency at $N_{\max} = 3$ (100\%) and $N_{\max} = 5$ (96.7\%), while NSGA-II shows the highest HV at $N_{\max} = 6$ (0.3609 vs.\ 0.3602 for SPEA2) but with 2.5$\times$ higher variance (std 0.015 vs.\ 0.006). This indicates that SPEA2 is more robust for deployment planning (consistent solutions), while NSGA-II explores more broadly but with less reproducibility.

\subsection{Layout Consistency Across Algorithms}

Table~\ref{tab:layout_all} presents the mode station patterns and their frequencies.

\begin{table}[htbp]
\caption{Mode Station Patterns and Layout Consistency}
\label{tab:layout_all}
\centering
\begin{tabular}{@{}llcc@{}}
\toprule
\textbf{$N_{\max}$} & \textbf{Algorithm} & \textbf{Mode Pattern} & \textbf{\%} \\
\midrule
\multirow{3}{*}{3} & NSGA-II & [1 0 0 1 1 1 1 0 0 1] & 80.0 \\
 & MOPSO & [1 0 0 1 1 1 1 0 0 1] & 96.7 \\
 & SPEA2 & [1 0 0 1 1 1 1 0 0 1] & 100.0 \\
\midrule
\multirow{3}{*}{4} & NSGA-II & [1 0 0 1 1 1 0 0 0 1] & 96.7 \\
 & MOPSO & [1 0 0 1 1 1 0 0 0 1] & 93.3 \\
 & SPEA2 & [1 0 0 1 1 1 0 0 0 1] & 90.0 \\
\midrule
\multirow{3}{*}{5} & NSGA-II & [1 0 0 1 1 0 0 0 0 1] & 86.7 \\
 & MOPSO & [1 0 0 1 1 0 0 0 0 1] & 70.0 \\
 & SPEA2 & [1 0 0 1 1 0 0 0 0 1] & 96.7 \\
\midrule
\multirow{3}{*}{6} & NSGA-II & [1 0 0 0 1 1 0 0 0 1] & 36.7 \\
 & MOPSO & [1 0 0 0 1 1 0 0 0 1] & 33.3 \\
 & SPEA2 & [1 0 0 0 1 1 0 0 0 1] & 46.7 \\
\bottomrule
\end{tabular}
\end{table}

\textbf{Why patterns shift:} The dominant pattern at $N_{\max} = 3$ opens stations at nodes 1, 4, 5, 6, 7, 10---covering corridor endpoints and the steep mid-route gradient zone (segments 4--8, gradients 1.7--1.8\%). As $N_{\max}$ increases, node 7 (at 69.7~km) closes first because it sits between two high-demand nodes (5 and 6), making its coverage redundant when per-node capacity increases. At $N_{\max} = 6$, the dominant pattern opens only nodes 1, 5, 6, 10, with node 5 absorbing the demand that previously required nodes 4 and 7.

\textbf{Why diversity increases at $N_{\max} = 6$:} With relaxed capacity constraints, multiple near-optimal configurations become feasible (e.g., opening nodes 1, 5, 6, 10 vs.\ nodes 1, 4, 5, 10), leading to lower mode frequency. This diversity provides decision-makers with flexible deployment options rather than indicating algorithmic weakness.

\subsection{Charger-Rate Sensitivity Analysis}

Table~\ref{tab:charger_rate} compares SPEA2 at 30~kW and 60~kW.

\begin{table}[htbp]
\caption{Charger-Rate Sensitivity (SPEA2, 30 Runs)}
\label{tab:charger_rate}
\centering
\begin{tabular}{@{}lccccc@{}}
\toprule
\textbf{Rate} & \textbf{$N_{\max}$} & \textbf{Open} & \textbf{Cost} & \textbf{HV} & \textbf{Layout} \\
\textbf{(kW)} & & \textbf{Stn} & \textbf{($\times 10^5$)} & \textbf{Norm} & \textbf{\%} \\
\midrule
\multirow{4}{*}{30} & 3 & 6.00 & 114.13 & 0.1109 & 100.0 \\
 & 4 & 5.00 & 103.36 & 0.2025 & 90.0 \\
 & 5 & 4.03 & 88.96 & 0.3170 & 96.7 \\
 & 6 & 3.87 & 87.85 & 0.3602 & 46.7 \\
\midrule
\multirow{4}{*}{60} & 3 & 3.00 & 95.52 & 0.2734 & 96.7 \\
 & 4 & 3.00 & 87.37 & 0.3734 & 100.0 \\
 & 5 & 2.00 & 72.30 & 0.4522 & 100.0 \\
 & 6 & 2.00 & 72.30 & 0.4798 & 100.0 \\
\bottomrule
\end{tabular}
\end{table}

\textbf{Why 60~kW enables further consolidation:} At 60~kW, each charger delivers twice the energy per unit time, effectively doubling the service rate $\mu$. This reduces the number of chargers needed per station (from 15 to 7 at $N_{\max} = 6$), allowing the two endpoint stations (nodes 1 and 10) to absorb all corridor demand. The extreme consolidation to 2 stations is physically justified by the corridor geometry: with 125~km total length, two well-placed stations at endpoints can serve all EVs within acceptable detour distances ($\leq$62.5~km).

\textbf{Why HV improves at 60~kW:} The 33.2\% HV improvement (0.3602 to 0.4798) reflects the expanded feasible region: with fewer stations and lower cost, the Pareto front extends further toward the ideal point in all three objectives simultaneously.

\subsection{Cost Sensitivity Analysis}
\label{sec:cost_sensitivity}

The perturbation study across 100 runs with varied cost multipliers (site cost: 0.84--1.17$\times$; charger cost: 0.84--1.18$\times$) and fleet sizes (3--7 EVs per fleet) demonstrates that the optimal station count is robust to cost variations: the mean open stations across 100 perturbed runs is 2.77 $\pm$ 0.43, with 100\% service coverage maintained. This indicates that the identified deployment patterns are not artifacts of specific cost assumptions but represent structurally optimal configurations for the corridor geometry.

The cost perturbation also reveals that the cost objective is less sensitive to station count than to charger unit price: a 10\% increase in charger cost increases total cost by 7.2\%, while reducing station count from 3 to 2 decreases cost by only 4.1\%. This suggests that charger procurement strategy has greater economic impact than station placement optimization.

\subsection{Convergence Behavior}

All algorithms converge within 500 generations. At $N_{\max} = 3$, convergence occurs by generation 50--100 (the solution space is constrained, limiting exploration). At $N_{\max} = 6$, convergence is slower (generations 100--200), reflecting the larger combinatorial space. Post-convergence HV stability is $\pm$1.5\% (std), confirming algorithmic robustness.

\subsection{Statistical Significance}

Wilcoxon rank-sum tests ($\alpha = 0.05$) for NSGA-II pairwise HV comparisons yield $p < 0.001$ for all adjacent $N_{\max}$ levels (3 vs.\ 4, 4 vs.\ 5, 5 vs.\ 6), as well as for the 30~kW vs.\ 60~kW comparison at $N_{\max} = 6$. All performance differences are statistically significant, confirming that the observed improvements are not due to random variation.

\subsection{Practical Deployment Recommendations}

\textbf{Budget-constrained:} $N_{\max} = 6$ at 30~kW, 4 stations (nodes 1, 5, 6, 10; 15 chargers; cost 87.49~$\times 10^5$~INR). Saves 23.2\% over 6-station baseline.

\textbf{Service-prioritized:} $N_{\max} = 4$ at 30~kW, 5 stations (nodes 1, 4, 5, 6, 10; 17 chargers; cost 103.36~$\times 10^5$~INR). Lowest waiting time (7.03~min) with 96.7\% consistency.

\textbf{High-power:} $N_{\max} \geq 5$ at 60~kW, 2 stations (nodes 1, 10; 7 chargers; cost 72.30~$\times 10^5$~INR). Saves 37.4\% over 6-station 30~kW baseline with 100\% consistency.

\subsection{Limitations and Comparison with Classical Approaches}

This study focuses on comparing metaheuristic algorithms rather than benchmarking against classical planning methods (p-median, greedy coverage, GIS-based heuristics) for several reasons: (i)~the multi-objective formulation with three conflicting objectives and mixed-integer decision variables (binary station + integer charger allocation) exceeds the scope of single-objective classical methods; (ii)~existing classical EVCS models \cite{13,14} do not incorporate terrain-aware energy models; and (iii)~the 100\% coverage constraint combined with the gradient-dependent demand model creates a problem structure where metaheuristics are necessary to explore the Pareto front. Future work will incorporate p-median and coverage heuristics as baseline comparisons.

The demand model assumes deterministic fleet operations (45 EVs, fixed schedules). This simplification is appropriate for the comparison study (all algorithms operate under identical conditions) but limits direct applicability to real-world stochastic demand. Seasonal tourism peaks and weather-driven demand variations remain important extensions.

\section{Conclusion}
\label{sec:conclusion}

This paper presented a comparative study of NSGA-II, MOPSO, and SPEA2 for multi-objective EVCS optimization on the mountainous Bagdogra--Gangtok corridor. Key findings:

\begin{enumerate}
\item Terrain-aware energy modeling increases per-trip energy consumption by 21.7\% over flat-road assumptions, directly influencing optimal station placement toward gradient transition points (nodes 5--6).

\item Relaxing $N_{\max}$ from 3 to 6 reduces stations by 36\%, cost by 23.2\%, and improves HV by 229\%, with all differences statistically significant ($p < 0.001$).

\item SPEA2 achieves the highest layout consistency (100\% at $N_{\max} = 3$ and $N_{\max} = 5$) and lowest HV variance (std 0.006), making it the most robust algorithm for deployment planning.

\item Doubling charger power from 30~kW to 60~kW enables station consolidation from 4 to 2 stations, reducing cost by 17.6\% with 100\% layout consistency.

\item Cost sensitivity analysis confirms that deployment patterns are robust to cost variations, with charger unit price having greater economic impact than station count.
\end{enumerate}

Future work will incorporate stochastic demand, seasonal patterns, grid constraints, and benchmark comparisons against classical planning methods (p-median, greedy coverage).

\section*{Acknowledgment}

The authors thank the School of Electrical Engineering, MIT MAHE Manipal for computational resources and technical support.

\begin{thebibliography}{00}
\bibitem{1} N. Tilly, T. Yigitcanlar, K. Degirmenci, and A. Paz, ``How sustainable is electric vehicle adoption? Insights from a PRISMA review,'' \textit{Sustainable Cities and Society}, vol. 117, p. 105950, 2024.
\bibitem{2} https://ev.tatamotors.com/tigor/ev.html
\bibitem{3} J. Mikolaj and L. Remek, ``Sustainable Maintenance of Low Level Road Network,'' \textit{MATEC Web of Conferences}, vol. 196, p. 04063, 2018.
\bibitem{4} I. Nutkani, H. Toole, N. Fernando, and L. P. C. Andrew, ``Impact of EV charging on electrical distribution network and mitigating solutions -- a review,'' \textit{IET Smart Grid}, 2024.
\bibitem{5} J. Dong, C. Liu, and Z. Lin, ``Charging infrastructure planning for promoting battery electric vehicles: An activity-based approach using multiday travel data,'' \textit{Transportation Research Part C}, vol. 38, pp. 44--55, 2014.
\bibitem{6} International Energy Agency, \textit{Global EV Outlook 2024}. IEA, 2024.
\bibitem{7} P. Joshi and C. Gokhale-Welch, ``Fundamentals of Electric Vehicles (EVs),'' National Renewable Energy Laboratory, pp. 1--30, 2022.
\bibitem{8} M. S. Mastoi et al., ``An In-Depth Analysis of Electric Vehicle Charging Station Infrastructure, Policy Implications, and Future Trends,'' \textit{Energy Reports}, vol. 8, pp. 11504--11529, 2022.
\bibitem{9} K. N. Hasan et al., ``Distribution Network Voltage Analysis with Data-Driven Electric Vehicle Load Profiles,'' \textit{Sustainable Energy, Grids and Networks}, vol. 36, 2023.
\bibitem{10} M. Aljaidi, N. Aslam, and O. Kaiwartya, ``Optimal Placement and Capacity of Electric Vehicle Charging Stations in Urban Areas: Survey and Open Challenges,'' \textit{IEEE Jordan JEEIT}, pp. 238--243, 2019.
\bibitem{11} N. Liu et al., ``A Charging Station Planning Model Considering the Charging Behavior of Private Car Users,'' \textit{10th RPGC}, pp. 225--231, 2021.
\bibitem{12} I. Chandra, N. K. Singh, and P. Samuel, ``A comprehensive review on coordinated charging of electric vehicles in distribution networks,'' \textit{Journal of Energy Storage}, vol. 89, p. 111659, 2024.
\bibitem{13} J. Banegas and J. Mamkhezri, ``A systematic review of GIS-based methods and criteria used for EVCS site selection,'' \textit{Environmental Science and Pollution Research}, vol. 30, no. 26, pp. 68054--68083, 2023.
\bibitem{14} B. Harshil and G. Nagababu, ``Strategies and models for optimal EV charging station site selection,'' \textit{IOP Conf. Series: Earth and Environmental Science}, vol. 1372, p. 012106, 2024.
\bibitem{15} S. Carpitella, X. Jia, and R. Narimani, ``A review on current trends on optimal EVCS location,'' \textit{Proc. 12th SESDE}, p. 003, 2024.
\bibitem{16} M. Y. Yang et al., ``Systematic Review of Social Equity for Installing Public EVCS,'' \textit{Construction Research Congress}, pp. 787--794, 2024.
\bibitem{17} S. Kumar et al., ``A comprehensive review of optimal placement of EVCS and sizing in active distribution network,'' \textit{Australian J. Electrical and Electronic Engineering}, 2024.
\bibitem{18} F. Ahmad et al., ``Optimal location of electric vehicle charging station and its impact on distribution network: A review,'' \textit{Energy Reports}, vol. 8, pp. 2314--2333, 2022.
\bibitem{19} M. Aljaidi et al., ``EV Charging Station Placement and Sizing Techniques: Survey, Challenges and Directions,'' \textit{ACIT}, pp. 1--6, 2022.
\bibitem{20} S. Islam et al., ``Effect of battery storage based EV chargers on harmonic profile of power system network,'' \textit{Journal of Energy Storage}, vol. 131, p. 117563, 2025.
\bibitem{21} M. Abdel-Basset et al., ``Metaheuristic Algorithms: A Comprehensive Review,'' in \textit{Computational Intelligence for Multimedia Big Data}, Academic Press, 2018, pp. 185-231.
\bibitem{22} M. Sharifi et al., ``Using NSGA II Algorithm for a Three-Objective Redundancy Allocation Problem,'' \textit{J. Optimization in Industrial Engineering}, vol. 19, pp. 87-95, 2015.
\bibitem{23} K. Deb et al., ``A fast and elitist multiobjective genetic algorithm: NSGA-II,'' \textit{IEEE Trans. Evolutionary Computation}, vol. 6, no. 2, pp. 182--197, 2002.
\bibitem{24} R. G. L. D'Souza et al., ``Improved NSGA-II based on a novel ranking scheme,'' \textit{Journal of Computing}, vol. 2, no. 2, pp. 91--95, 2010.
\bibitem{25} A. Prakash P. and R. Radha, ``Optimization of electric vehicle routing: A multi-criteria intelligent decision framework for efficient and scalable mobility,'' \textit{Energy Reports}, vol. 15, Jun. 2026, Art. no. 108912.
\end{thebibliography}

\end{document}
